% Итоговый отчёт по тестовому заданию. Начальное состояние памяти из энкодера RMT для GradMem.
% Компилируется pdfLaTeX (Overleaf или TinyTeX) без дополнительных пакетов.
\documentclass[a4paper,11pt]{article}

\usepackage[T2A]{fontenc}
\usepackage[utf8]{inputenc}
\usepackage[russian]{babel}
\usepackage[top=1.6cm,left=2.2cm,right=2.2cm,bottom=1.8cm]{geometry}
\usepackage{amsmath,amssymb}
\usepackage{booktabs}
\usepackage{graphicx}
\usepackage{array}
\usepackage{microtype}
\usepackage[font=small,labelsep=period]{caption}
\usepackage[colorlinks=true,linkcolor=black,citecolor=black,urlcolor=black]{hyperref}

\setlength{\parindent}{1.2em}
\setlength{\parskip}{2pt}
\setlength{\textfloatsep}{8pt plus 2pt minus 2pt}
\setlength{\floatsep}{8pt plus 2pt minus 2pt}
\setlength{\intextsep}{8pt plus 2pt minus 2pt}
\setlength{\abovecaptionskip}{3pt}
\setlength{\tabcolsep}{4pt}
\newcommand{\pmm}{\ensuremath{\,\pm\,}}
\renewcommand{\arraystretch}{1.08}
\newcommand{\cm}{\mathcal{L}}

\title{Начальное состояние памяти из энкодера RMT для GradMem\\[2pt]\large Отчёт по тестовому заданию}
\author{Александр Семин}
\date{}

\begin{document}
\maketitle
\vspace{-1.2em}

\begin{abstract}
Исследуется, может ли начальное состояние памяти, вычисленное энкодером RMT за один прямой проход, заменить часть градиентных шагов записи в память GradMem. На задаче ассоциативного поиска MQAR с 16 и 32 парами из ключа и значения комбинированная модель с $N$ шагами записи сравнивается с GradMem с $N+1$ и $N+2$ шагами и с RMT с $N+1$ обновлениями памяти. Относительно GradMem, инициализированного весами RMT, как предусмотрено заданием, энкодер заменяет один шаг записи при 16 парах и не менее двух шагов при 32 парах. Относительно GradMem, обученного с нуля, замена неполная. При 16 парах комбинированная модель с одним шагом уступает GradMem с двумя шагами 3,0 пункта точного совпадения и работает на 28\,\% быстрее, а при 32 парах модель с двумя шагами уступает GradMem с тремя шагами 6,4 пункта. Дополнительные проходы RMT выигрыша не дают. Повторное использование прямого прохода шага записи даёт выигрыш, когда его результат добавляется к градиентному шагу.
\end{abstract}

\section{Постановка задачи}
Задание состоит из трёх частей.
\begin{enumerate}\setlength{\itemsep}{0pt}
\item Сравнить RMT, GradMem, комбинацию энкодера RMT с GradMem при $N$ обновлениях памяти и RMT с $N+1$ обновлениями.
\item Проверить, существует ли постановка, в которой начальное состояние от энкодера позволяет уменьшить число обратных проходов записи в память GradMem.
\item Дополнительно проверить, можно ли переиспользовать вычисления прямого прохода операции записи GradMem для обновления памяти, сохранив градиентный шаг.
\end{enumerate}

Память в обоих методах состоит из $m=8$ векторов размерности $d=128$. RMT \cite{rmt} записывает контекст $x$ прямым проходом трансформера по последовательности, в которой память стоит до и после контекста, и берёт новые векторы памяти с последних $m$ позиций
\begin{equation}
M' = \operatorname{Tr}_\theta\big([\,M;\,x;\,M\,]\big)_{\text{последние } m}. \label{eq:rmt}
\end{equation}
RMT с $N$ обновлениями повторяет этот проход $N$ раз. GradMem \cite{gradmem} записывает контекст $K$ градиентными шагами по функции потерь восстановления контекста при замороженных весах
\begin{equation}
M_{k+1} = M_k - \alpha\,\nabla_M \cm_{\text{write}}(M_k;\theta,x), \qquad \cm_{\text{write}} = -\frac{1}{|x|}\sum_{t}\log p_\theta\big(x_t \mid [\,M;\,x_{<t}\,]\big), \label{eq:gm}
\end{equation}
где $\alpha$ обозначает внутренний шаг обучения. Чтение выполняется по $[\,M_K;\,q\,]$, где $q$ содержит запросы, а внешняя функция потерь обучает веса через все шаги записи с производными второго порядка. Каждый шаг записи требует прямого и обратного прохода, поэтому сокращение $K$ напрямую уменьшает стоимость записи.

Комбинированная модель повторяет идею полуамортизированного вывода \cite{savae}, в которой кодировщик даёт начальное приближение, а несколько градиентных шагов его уточняют. Энкодер $E_\phi$, копия обученного RMT, записывает контекст за один проход \eqref{eq:rmt}, после чего GradMem выполняет $K$ шагов \eqref{eq:gm} от этой памяти. Обе части обучаются совместно и используют общую таблицу эмбеддингов.

\section{Экспериментальная постановка}
\paragraph{Задача и модель.} Используется задача MQAR из набора Zoology \cite{zoology}. Контекст содержит $n$ пар из ключа и значения, взятых из словаря в 8192 токена, затем идут $n$ запросов к ключам в случайном порядке, и модель должна назвать значение каждого ключа. Основные опыты проводятся при $n=16$, проверка на более трудной задаче при $n=32$. Все модели имеют одинаковый размер. Это трансформер Llama с 4 слоями, размерностью 128 и 4 головами внимания. Обучающая выборка содержит $10^6$ примеров, валидационная 3000, а итоговые числа получены на отдельной тестовой выборке из 3000 примеров с другим зерном генератора. Обучение идёт с размером батча 64, скоростью обучения $10^{-4}$ с прогревом на 1000 шагов, оценкой каждые 500 шагов и не более чем 60 тыс. шагов. Обучение прекращается досрочно, если точное совпадение на валидации не ниже 99\,\% в трёх оценках подряд. Обучение проводилось на GPU Tesla T4 и GTX~1080~Ti, а скорость всех моделей измерена на одной GTX~1080~Ti при размере батча 64.

\paragraph{Метрики.} Точность по токенам (token accuracy) равна доле правильно названных значений. Точное совпадение (exact match, EM) равно доле примеров, в которых правильны все $n$ ответов. При независимых ошибках с точностью по токенам $p$ получаем $\mathrm{EM}\approx p^{\,n}$, поэтому при $n=32$ даже 90\,\% по токенам дают лишь около 3\,\% точного совпадения.
\paragraph{Обозначения.} E1 обозначает RMT с одним обновлением, обученный с функцией потерь восстановления. RMT$\times K$ обозначает дообучение E1 с $K$ проходами записи, в таблицах это RMT с соответствующим $K$. GradMem от E1 начинает с весов E1, как требует схема задания. GradMem с нуля начинает со случайных весов и служит более сильной контрольной моделью. Комбинированная модель обозначена как RMT+GradMem.

\section{Подготовительные результаты}
\paragraph{Функция потерь восстановления для RMT.} При обучении только по функции потерь чтения RMT при 4, 8 и 16 парах оставался ровно на уровне $1/n$. Модель выучивала ответ любым значением из контекста и не хранила связь ключа со значением. К функции потерь RMT была добавлена та же функция восстановления контекста, что используется в GradMem \cite{gradmem}, с весом 1. После этого E1 при 8 парах достиг 99,5\,\% по токенам и 96,2\,\% точного совпадения, а проверка с перестановкой значений между ключами подтвердила, что модель хранит именно связи. При 16 парах E1 правильно читает пары с первой по пятую (97,0\,\%), а остальные лишь на уровне 8,4\,\%, хотя функция потерь восстановления, равная 0,276, показывает, что в памяти хранится почти весь контекст.

\paragraph{Масштаб начальной памяти.} Выход последнего слоя проходит через RMSNorm \cite{rmsnorm} с обучаемым коэффициентом $\gamma$, поэтому норма памяти энкодера равна $\sqrt{md}\,\mathrm{RMS}(\gamma)$. Для E1 это около 102, тогда как память GradMem имеет норму от 2 до 10. Если функция потерь инвариантна к масштабу входа, то $\nabla\cm(tM)=\nabla\cm(M)/t$ и $\langle\nabla\cm(M),M\rangle=0$ \cite{scaleinv}. Для архитектуры с нормализацией перед слоем \cite{prenorm} это выполняется приближённо при большой норме памяти, когда она превосходит вклад слоёв в остаточный поток. Тогда относительный шаг записи убывает как
\begin{equation}
\frac{\lVert M_{k+1}-M_k\rVert}{\lVert M_k\rVert} = \frac{\alpha\,\lVert\nabla\cm(M_k)\rVert}{\lVert M_k\rVert} \propto \frac{\alpha}{\lVert M_k\rVert^{2}}, \label{eq:relstep}
\end{equation}
и при норме 102 шаг меняет память лишь на тысячные доли процента. Действительно, комбинированная модель с исходной памятью энкодера при всех шагах $\alpha\in\{0{,}1;0{,}2;0{,}4\}$ и $K\in\{1,2\}$ оставалась на уровне E1, то есть от 36,8 до 37,5\,\% по токенам. Поэтому начальная память масштабируется к фиксированной длине
\begin{equation}
M_0 = \ell\,\frac{E_\phi(x)}{\lVert E_\phi(x)\rVert}, \qquad \ell = 5. \label{eq:scaled}
\end{equation}
Длина $\ell$ выбрана по пробным запускам на 10 тыс. шагов. При $\ell=5$ точное совпадение составило 98,5\,\%, при $\ell=10$ оно равно 96,4\,\%, при $\ell=20$ равно 77,1\,\%. Случайное направление той же длины вместо выхода энкодера дало 0,1\,\%, то есть полезен именно выход энкодера.

\paragraph{Подбор шага записи.} Для GradMem от E1 и для комбинированной модели шаг $\alpha$ подобран на прогонах по 10 тыс. шагов по одной сетке значений $\{0{,}1;0{,}2;0{,}4\}$. Для GradMem от E1 выбраны $\alpha=0{,}4$ при $K=1$ и $\alpha=0{,}1$ при $K=2$ и $K=3$, для комбинированной модели $\alpha=0{,}1$ при $K=1$ и $\alpha=0{,}4$ при $K=2$. Для GradMem с нуля при 16 парах взят исходный шаг метода $\alpha=0{,}4$ при $K\le2$ и $\alpha=0{,}1$ при $K=3$. При 32 парах для обоих методов использовано $\alpha=0{,}1$, кроме комбинированной модели с $K=2$.

\section{Сравнение с большим числом обновлений памяти}
Комбинированная модель с $N$ шагами записи сравнивается с GradMem с $N+1$ и $N+2$ шагами и с RMT с $N+1$ обновлениями. Каждый блок таблиц~\ref{tab:main} и~\ref{tab:p32} начинается с комбинированной модели, за которой идут модели, с которыми она сравнивается.

\paragraph{Как читать таблицы.} Столбец $K$ равен числу обновлений памяти. Для GradMem это градиентные шаги записи, для RMT+GradMem градиентные шаги после одного прохода энкодера, для RMT прямые проходы записи. Exact match и token accuracy даны в процентах на тестовой выборке. Запись $96{,}2\pm0{,}7$ означает среднее по трём зёрнам и стандартное отклонение между ними, а для строк с одним зерном разброс не указан. Время на батч равно числу миллисекунд, за которые модель записывает в память и читает из неё батч из 64 примеров на GTX~1080~Ti, поэтому меньшее время означает более быструю модель. Столбец $\Delta$ показывает, на сколько пунктов exact match строки выше (плюс) или ниже (минус), чем у RMT+GradMem того же блока.

\begin{table}[t]
\centering\footnotesize
\caption{Задача с 16 парами. Столбцы описаны в начале раздела. Знак $\ast$ означает обучение до 120 тыс. шагов.}
\label{tab:main}
\resizebox{\textwidth}{!}{\begin{tabular}{l c c c c c c}
\toprule
Метод & $K$ & Зёрна & Exact match, \% & Token accuracy, \% & Время на батч, мс & $\Delta$ \\
\midrule
RMT (E1) & 1 & 1 & 0,0 & 36,1 & 12,3 & \\
\midrule
\multicolumn{7}{l}{\textit{Блок $N=1$}} \\
\textbf{RMT+GradMem}$^\ast$ & 1 & 3 & \textbf{96,2\pmm0,7} & 99,8\pmm0,0 & \textbf{31,9} & \\
GradMem с нуля & 2 & 3 & 99,2\pmm0,1 & 99,9\pmm0,0 & 44,1 & $+3{,}0$ \\
GradMem от E1$^\ast$ & 2 & 3 & 98,5\pmm0,8 & 99,9\pmm0,0 & 47,5 & $+2{,}3$ \\
GradMem с нуля & 3 & 3 & 99,3\pmm0,4 & 100,0\pmm0,0 & 67,7 & $+3{,}1$ \\
GradMem от E1 & 3 & 3 & 98,7\pmm0,6 & 99,9\pmm0,0 & 61,3 & $+2{,}5$ \\
RMT & 2 & 1 & 0,0 & 36,9 & 18,7 & $-96{,}2$ \\
\midrule
\multicolumn{7}{l}{\textit{Блок $N=2$}} \\
\textbf{RMT+GradMem} & 2 & 3 & \textbf{99,0\pmm0,2} & 99,9\pmm0,0 & \textbf{52,0} & \\
GradMem с нуля & 3 & 3 & 99,3\pmm0,4 & 100,0\pmm0,0 & 67,7 & $+0{,}3$ \\
GradMem от E1 & 3 & 3 & 98,7\pmm0,6 & 99,9\pmm0,0 & 61,3 & $-0{,}3$ \\
RMT & 3 & 1 & 0,0 & 36,9 & 24,7 & $-99{,}0$ \\
\bottomrule
\end{tabular}}
\end{table}

\paragraph{Задача с 16 парами.} Таблица~\ref{tab:main} отвечает на первые две части задания. RMT$\times2$ и RMT$\times3$ остаются на уровне E1, около 37\,\% по токенам и 0\,\% точного совпадения, и по-прежнему читают только первые пять пар. Следовательно, дополнительные прямые проходы RMT не заменяют градиентных шагов. Относительно GradMem от E1 комбинированная модель с двумя шагами не уступает GradMem с тремя шагами (99,0 против 98,7\,\%) и работает на 15\,\% быстрее, то есть энкодер заменяет один шаг записи. Комбинированная модель с одним шагом уступает GradMem от E1 с двумя шагами 2,3 пункта. Относительно GradMem с нуля замена одного шага стоит 3,0 пункта точного совпадения и сокращает время обработки на 28\,\%, а замена двух шагов стоит 3,1 пункта и сокращает время обработки на 53\,\%. Комбинированная модель с двумя шагами отстаёт от GradMem с нуля с тремя шагами всего на 0,3 пункта при сокращении времени обработки на 23\,\% и выходит на 99\,\% точного совпадения за 13,0 тыс. шагов обучения, примерно как GradMem с нуля с тремя шагами (12,3 тыс.).

Комбинированная модель имеет 3,15 млн параметров против 2,10 млн у GradMem, так как энкодер является второй копией E1 с общей таблицей эмбеддингов. Проход энкодера стоит около 5 мс на батч, а каждый градиентный шаг записи от 14 до 20 мс, то есть энкодер обходится примерно в четверть или треть шага.

\begin{table}[t]
\centering\footnotesize
\caption{Задача с 32 парами. Обозначения как в таблице~\ref{tab:main}. GradMem от E1 обучен с шагом $\alpha=0{,}1$, если шаг не указан.}
\label{tab:p32}
\resizebox{\textwidth}{!}{\begin{tabular}{l c c c c c c}
\toprule
Метод & $K$ & Зёрна & Exact match, \% & Token accuracy, \% & Время на батч, мс & $\Delta$ \\
\midrule
RMT (E1 при 32 парах) & 1 & 1 & 0,0 & 5,0 & 12,9 & \\
\midrule
\multicolumn{7}{l}{\textit{Блок $N=1$}} \\
RMT+GradMem & 1 & 1 & 2,4 & 89,9 & 37,7 & \\
GradMem с нуля & 2 & 3 & 91,1\pmm1,6 & 99,7\pmm0,1 & 57,1 & $+88{,}7$ \\
GradMem от E1 & 2 & 1 & 0,0 & 62,1 & 55,1 & $-2{,}4$ \\
GradMem с нуля & 3 & 3 & 99,1\pmm0,2 & 100,0\pmm0,0 & 80,4 & $+96{,}7$ \\
GradMem от E1 & 3 & 1 & 9,5 & 92,3 & 80,7 & $+7{,}1$ \\
RMT & 2 & 1 & 0,0 & 15,8 & 20,6 & $-2{,}4$ \\
\midrule
\multicolumn{7}{l}{\textit{Блок $N=2$}} \\
\textbf{RMT+GradMem} & 2 & 3 & \textbf{92,7\pmm0,2} & 99,8\pmm0,0 & \textbf{67,0} & \\
GradMem с нуля & 3 & 3 & 99,1\pmm0,2 & 100,0\pmm0,0 & 80,4 & $+6{,}4$ \\
GradMem от E1 & 3 & 1 & 9,5 & 92,3 & 80,7 & $-83{,}2$ \\
GradMem с нуля & 4 & 3 & 99,3\pmm0,1 & 100,0\pmm0,0 & 110,9 & $+6{,}6$ \\
GradMem от E1 & 4 & 3 & 21,6\pmm11,4 & 93,7\pmm2,6 & 107,6 & $-71{,}1$ \\
GradMem от E1, $\alpha=0{,}05$ & 4 & 1 & 66,8 & 98,6 & 104,6 & $-25{,}9$ \\
RMT & 3 & 1 & 0,0 & 12,9 & 27,5 & $-92{,}7$ \\
\midrule
\multicolumn{7}{l}{\textit{Блок $N=3$}} \\
\textbf{RMT+GradMem} & 3 & 3 & \textbf{98,5\pmm0,4} & 100,0\pmm0,0 & \textbf{87,3} & \\
GradMem с нуля & 4 & 3 & 99,3\pmm0,1 & 100,0\pmm0,0 & 110,9 & $+0{,}8$ \\
GradMem от E1 & 4 & 3 & 21,6\pmm11,4 & 93,7\pmm2,6 & 107,6 & $-76{,}9$ \\
GradMem от E1, $\alpha=0{,}05$ & 4 & 1 & 66,8 & 98,6 & 104,6 & $-31{,}7$ \\
RMT & 4 & 1 & 0,0 & 12,6 & 34,0 & $-98{,}5$ \\
\bottomrule
\end{tabular}}
\end{table}

\paragraph{Задача с 32 парами.} Относительно GradMem от E1 комбинированная модель с двумя шагами (92,7\,\%) превосходит GradMem от E1 с четырьмя шагами (21,6\,\%, а с уменьшенным шагом $\alpha=0{,}05$ равно 66,8\,\%) и работает на 38\,\% быстрее, то есть экономия составляет не менее двух шагов записи. Относительно GradMem с нуля замена одного шага стоит больше, чем при 16 парах. Комбинированная модель с двумя шагами уступает GradMem с тремя шагами 6,4 пункта при сокращении времени обработки на 17\,\%, а GradMem с четырьмя шагами 6,6 пункта при сокращении времени обработки на 40\,\%. С тремя шагами комбинированная модель отстаёт от GradMem с нуля с четырьмя шагами на 0,8 пункта и работает на 21\,\% быстрее. Одного шага записи при 32 парах недостаточно ни в одной постановке.

RMT с двумя, тремя и четырьмя проходами записи при 32 парах не даёт ни одного полностью правильного примера. Точность по токенам растёт с 5,0\,\% у E1 до 15,8\,\% при двух проходах и снижается до 12,9 и 12,6\,\% при трёх и четырёх. Первые пять пар читаются с точностью от 32 до 39\,\%, остальные от 9 до 12\,\%. Таким образом, и при 32 парах дополнительные прямые проходы RMT не заменяют градиентных шагов, а каждый следующий проход только увеличивает время обработки.

Некоторые значения при 32 парах выглядят выбросами. Комбинированная модель с одним шагом и GradMem от E1 с тремя и четырьмя шагами отвечают правильно на 90 до 94\,\% запросов, а низкое точное совпадение возникает из соотношения $\mathrm{EM}\approx p^{32}$. Например, $0{,}899^{32}\approx 3{,}3\,\%$ при измеренных 2,4\,\%. GradMem от E1 при 32 парах страдает от трёх причин. Чтение E1 при 32 парах не обучилось (5,0\,\% по токенам), шаг $\alpha=0{,}1$ слишком велик, а за 60 тыс. шагов обучение не завершилось. Поэтому основной контрольной моделью при 32 парах служит GradMem с нуля.

\section{Абляции и контрольные опыты}
Абляции проведены на одном зерне с бюджетом 60 тыс. шагов. Для сравнения при $K=1$ используется основная комбинированная модель на 60 тыс. шагов с тем же зерном, у которой точное совпадение на тесте равно 91,6\,\%.

\begin{table}[t]
\centering\footnotesize
\caption{Абляции комбинированной модели, 16 пар, одно зерно. Exact match на тесте в процентах. Перемешанный старт означает, что каждый пример получает память энкодера от другого примера.}
\label{tab:abl}
\resizebox{\textwidth}{!}{\begin{tabular}{l c c c c c}
\toprule
Вариант & $K$ & Exact match & Exact match, перемешанный старт & Exact match, 1 шаг при тесте & 99\,\% exact match, шаг \\
\midrule
Основная модель, 60 тыс. шагов & 1 & 91,6 & 48,9 & 91,6 & нет \\
С проекцией начальной памяти & 1 & 95,3 & 71,0 & 95,3 & нет \\
С функцией восстановления для энкодера & 1 & 87,7 & 14,6 & 87,7 & нет \\
С замороженным энкодером & 1 & 76,0 & 0,0 & 76,0 & нет \\
\midrule
Основная модель & 2 & 98,8 & 98,8 & 33,4 & 11,5 тыс. \\
С функцией восстановления для энкодера & 2 & 98,9 & 98,9 & 39,2 & 16,0 тыс. \\
С замороженным энкодером & 2 & 99,0 & 95,6 & 25,2 & 48,0 тыс. \\
\bottomrule
\end{tabular}}
\end{table}

Таблица~\ref{tab:abl} даёт три вывода. Во-первых, энкодер необходимо дообучать вместе с GradMem. Замороженный энкодер теряет 15,6 пункта при $K=1$ и в четыре раза замедляет выход на 99\,\% при $K=2$. Во-вторых, обучаемая проекция начальной памяти улучшает результат при $K=1$ на 3,7 пункта, а дополнительная функция восстановления для энкодера ухудшает его. В-третьих, содержимое памяти энкодера используется тогда, когда шагов записи недостаточно. При $K=1$ перемешивание стартов между примерами снижает точное совпадение с 96,2 до $61{,}3\pmm12{,}0\,\%$ для трёх зёрен основной модели, а при $K=2$ почти не влияет (99,0 и 98,6\,\%). При 32 парах та же картина сдвигается на один шаг. Перемешивание снижает результат при $K=2$ с 92,7 до 86,5\,\%, а при $K=3$ лишь с 98,5 до 97,6\,\%. Модели, обученные с $K$ шагами, требуют ровно $K$ шагов при тестировании. Комбинированная модель с $K=1$, оценённая с двумя шагами, достигает 98,0\,\%, тогда как GradMem с $K=3$ при двух шагах падает до 63,5\,\%.

\section{Переиспользование прямого прохода записи}\label{sec:fwd}
Каждый шаг GradMem выполняет прямой проход по $[\,M_k;\,x\,]$, но из него используется только значение $\cm_{\text{write}}$. Проверены четыре способа использовать остальные вычисления. Все варианты обучены при 16 парах от E1 и с нуля. Последние три используют те же шаги $\alpha$, что соответствующий GradMem. Слияние обучено с шагами комбинированной модели ($\alpha=0{,}1$ при $K=1$ и $0{,}4$ при $K=2$), которые в трёх из четырёх постановок отличаются от шагов GradMem, но его отставание намного больше, чем можно объяснить этим различием.

\paragraph{Слияние.} Этот вариант предложен в задании. К контексту добавляются $m$ слотов памяти, как в RMT, и тот же прямой проход даёт новую память $F_\theta(M_k,x)$ с последних позиций. Из-за каузального внимания эти слоты не влияют на позиции контекста, поэтому функция записи и её градиент совпадают с GradMem, что проверено тестом. Обновление имеет вид
\begin{equation}
M_{k+1} = \ell\,\frac{F_\theta(M_k,x)}{\lVert F_\theta(M_k,x)\rVert} - \alpha\,\nabla_M\cm_{\text{write}}(M_k). \label{eq:fused}
\end{equation}

\paragraph{Добавление контекста.} Градиентный шаг сохраняется полностью, а к нему добавляется вклад скрытых состояний $h_t$ последнего слоя на позициях контекста, которые вычислены тем же проходом. Каждый слот памяти собирает их вниманием, аналогично записи в быстрые веса \cite{fastweights} и обучению во время вывода \cite{ttt}
\begin{equation}
M_{k+1} = M_k - \alpha\nabla_M\cm_{\text{write}}(M_k) + \sigma(w)\,\lVert M_k\rVert\frac{D}{\lVert D\rVert}, \quad D_j = \sum_t \operatorname{softmax}_t\!\Big(\tfrac{(W_q M_{k,j})^{\top} W_k h_t}{\sqrt d}\Big) W_v h_t. \label{eq:ctx}
\end{equation}
Затвор начинается с $\sigma(-4)\approx 0{,}018$, поэтому модель стартует как обычный GradMem. Добавляется 49 тыс. параметров, около 2\,\%.

\paragraph{Взвешивание токенов.} Потери отдельных токенов из прямого прохода задают веса $w_t \propto \operatorname{sg}(\cm_t^{\,\gamma})$, где $\cm_t$ обозначает потерю на токене $t$, а $\operatorname{sg}$ остановку градиента, при $\gamma=1$ со средним весом 1, так что градиент сосредоточен на плохо записанных парах, как в focal loss \cite{focal} и взвешивании токенов \cite{tokenweight}.

\paragraph{Шаг Поляка.} Значение функции потерь из прямого прохода задаёт размер шага для каждого примера $\alpha_k=\min\big(\alpha,\ \cm(M_k)/\lVert\nabla\cm(M_k)\rVert^2\big)$ \cite{polyak,sps}.

\begin{table}[t]
\centering\footnotesize
\caption{Переиспользование прямого прохода, 16 пар, exact match на тесте в процентах. Знаком $\ast$ отмечен GradMem от E1 при $K=2$, обученный до 120 тыс. шагов (на валидации после 60 тыс. шагов, как у остальных вариантов, 97,4\,\%). Знаком $\ddagger$ отмечена ячейка, где 90\,\% достигли два зерна из трёх. Столбец с шагом показывает, когда EM на валидации достигает 90\,\% при $K=2$ от E1. С нуля обычный GradMem и слияние обучены на трёх зёрнах, остальные варианты на одном.}
\label{tab:fwd}
\begin{tabular}{l c c c c c c}
\toprule
& \multicolumn{2}{c}{От E1, $K=1$} & \multicolumn{2}{c}{От E1, $K=2$} & \multicolumn{2}{c}{С нуля} \\
\cmidrule(lr){2-3}\cmidrule(lr){4-5}\cmidrule(lr){6-7}
Вариант & Зёрна & Exact match & Exact match & 90\,\% exact match, шаг & $K=1$ & $K=2$ \\
\midrule
GradMem & 3 & 33,0\pmm16,1 & 98,5$^\ast$ & 26,5 тыс. & 76,4\pmm6,3 & 99,2\pmm0,1 \\
Слияние & 3 & 8,9\pmm3,0 & 94,6\pmm1,6 & 40,2 тыс. & 28,4\pmm23,7 & 96,7\pmm0,8 \\
\textbf{Добавление контекста} & 3 & \textbf{59,2\pmm10,5} & \textbf{97,7\pmm0,5} & \textbf{19,2 тыс.} & 59,9 & 98,9 \\
Взвешивание токенов & 3 & 50,3\pmm12,6 & 93,0\pmm5,1 & 23,2 тыс.$^\ddagger$ & 66,5 & 99,2 \\
Шаг Поляка & 1 & 2,5 & 85,3 & нет & 10,6 & 98,9 \\
\bottomrule
\end{tabular}
\end{table}

Таблица~\ref{tab:fwd} показывает, что слияние в форме \eqref{eq:fused} ухудшает результат во всех постановках. При $K=2$ длина градиентного шага составляет лишь от 12 до 19\,\% длины нормированной записи RMT, поэтому память на каждом шаге в основном заменяется записью RMT, а запись RMT на этой задаче сама не обучается чтению (таблица~\ref{tab:main}). Добавление контекста по формуле \eqref{eq:ctx} даёт небольшой, но устойчивый выигрыш от E1. При $K=2$ точное совпадение на валидации после 60 тыс. шагов равно 98,1 против 97,4\,\% у GradMem, а 90\,\% достигаются на 27\,\% раньше. При $K=1$ результат на валидации выше, чем у GradMem, на каждом из трёх зёрен, а на тесте в среднем составляет 59,2 против 33,0\,\%, хотя разброс велик. Затвор во время обучения растёт с 0,018 примерно до 0,05, то есть модель действительно использует добавку. Время обработки батча не меняется (46,3 против 47,5 мс при $K=2$). Взвешивание токенов нестабильно, а шаг Поляка в начале обучения в 60 раз меньше обычного и тормозит обучение. При обучении с нуля ни один вариант не превосходит обычный GradMem, и ни один из них не позволяет заменить второй шаг записи первым.

\section{Выводы и ограничения}
\begin{enumerate}\setlength{\itemsep}{0pt}
\item Дополнительные прямые проходы RMT не улучшают чтение. При 16 парах RMT$\times2$ и RMT$\times3$ остаются на уровне E1, а при 32 парах RMT с двумя, тремя и четырьмя проходами даёт 0\,\% точного совпадения, поэтому такие проходы не заменяют градиентных шагов.
\item Относительно GradMem от E1, то есть в постановке задания, начальное состояние энкодера заменяет один шаг записи при 16 парах и не менее двух шагов при 32 парах и сокращает время обработки на 15 и 38\,\% соответственно.
\item Относительно GradMem с нуля замена неполная. При 16 парах замена одного шага стоит 3,0 пункта точного совпадения и сокращает время обработки на 28\,\%, а комбинированная модель с двумя шагами отстаёт от GradMem с тремя шагами на 0,3 пункта при сокращении времени обработки на 23\,\%. При 32 парах замена шага стоит 6,4 пункта при сокращении времени обработки на 17\,\%.
\item Энкодер должен обучаться вместе с GradMem, а начальную память нужно приводить к масштабу градиентных шагов. Содержимое этой памяти используется тогда, когда шагов записи недостаточно.
\item Прямой проход шага записи даёт выигрыш, когда его результат добавляется к градиентному шагу. Замена памяти этим результатом ухудшает запись.
\end{enumerate}

Ограничения таковы. Абляции, RMT$\times K$ и варианты с нуля из раздела~\ref{sec:fwd}, кроме слияния, проведены на одном зерне. Шаг $\alpha$ при 32 парах не подбирался отдельно. Все выводы получены на синтетической задаче MQAR с маленькой моделью и требуют проверки на естественных текстах.

\begin{thebibliography}{99}\small
\bibitem{gradmem} Y. Kuratov, M. Kairov, A. Bulatov, I. Rodkin, M. Burtsev. GradMem: Learning to Write Context into Memory with Test-Time Gradient Descent. ICML, 2026.
\bibitem{rmt} A. Bulatov, Y. Kuratov, M. Burtsev. Recurrent Memory Transformer. NeurIPS, 2022.
\bibitem{zoology} S. Arora, S. Eyuboglu, A. Timalsina, I. Johnson, M. Poli, J. Zou, A. Rudra, C. Ré. Zoology: Measuring and Improving Recall in Efficient Language Models. arXiv 2312.04927, 2023.
\bibitem{savae} Y. Kim, S. Wiseman, A. C. Miller, D. Sontag, A. M. Rush. Semi-Amortized Variational Autoencoders. ICML, 2018.
\bibitem{rmsnorm} B. Zhang, R. Sennrich. Root Mean Square Layer Normalization. NeurIPS, 2019.
\bibitem{scaleinv} S. Arora, Z. Li, K. Lyu. Theoretical Analysis of Auto Rate-Tuning by Batch Normalization. ICLR, 2019.
\bibitem{prenorm} R. Xiong et al. On Layer Normalization in the Transformer Architecture. ICML, 2020.
\bibitem{fastweights} I. Schlag, K. Irie, J. Schmidhuber. Linear Transformers Are Secretly Fast Weight Programmers. ICML, 2021.
\bibitem{ttt} Y. Sun et al. Learning to (Learn at Test Time): RNNs with Expressive Hidden States. arXiv 2407.04620, 2024.
\bibitem{focal} T.-Y. Lin, P. Goyal, R. Girshick, K. He, P. Dollár. Focal Loss for Dense Object Detection. ICCV, 2017.
\bibitem{tokenweight} F. Helm, N. Daheim, I. Gurevych. Token Weighting for Long-Range Language Modeling. Findings of NAACL, 2025.
\bibitem{polyak} B. T. Polyak. Introduction to Optimization. Optimization Software, 1987.
\bibitem{sps} N. Loizou, S. Vaswani, I. Laradji, S. Lacoste-Julien. Stochastic Polyak Step-size for SGD. AISTATS, 2021.
\end{thebibliography}

\end{document}
